\subsection{معماری و پیاده‌سازی}
کلاس \lr{LiLLaMACompressor} در مسیر \lr{neuralops/compression\_engine/lillama\_compressor.py} واقع شده و تحت شناسه‌های \lr{lillama}، \lr{lillama\_compressor}، \lr{lillama\_class} و \lr{lillama\_factory} در رجیستری فشرده‌سازی ثبت شده است. مؤلفه‌ها و مکانیزم‌های کلیدی پیاده‌سازی عبارتند از:
\begin{itemize}
    \item \textbf{مکانیزم قلاب‌های رهگیری پویا (\lr{Forward Activation Hooks}):} ماژول قلاب‌های موقتی را بر روی تمامی لایه‌های خطی هدف نصب نموده و در یک گذر پیشرو بر روی $N_{\text{batch}}$ دسته داده متنی، تانسورهای فعال‌سازی ورودی و خروجی را جمع‌آوری و برای مدیریت حافظه به پردازنده میزبان انتقال می‌دهد.
    \item \textbf{سازگاری کامل با لایه‌های \lr{Conv1D} و شبکه‌های دیکودر:} لایه‌های پیچشی ۱بعدی (نظیر ساختار \lr{c\_attn} در \lr{GPT-2}) از طریق ترانهاده‌سازی خودکار و بازآرایی ابعاد پشتیبانی می‌شوند.
    \item \textbf{فیلتر هوشمند لایه‌ها و لایه‌های بحرانی:} لایه‌های حساس نظیر سر رده‌بندی، لایه‌های نشاندن کلمات (\lr{Embeddings}) و لایه‌های با ابعاد کمتر از $\text{min\_layer\_dim}=8$ از فرآیند فشرده‌سازی مستثنی می‌شوند تا پایداری خروجی تضمین گردد.
    \item \textbf{تضمین عدم اجرای تکراری (\lr{Idempotency Guard}):} با تنظیم پرچم اختصاصی بر روی مدل پس از فشرده‌سازی، از اجرای ناخواسته فشرده‌سازی چندگانه جلوگیری می‌شود.
\end{itemize}
